\section{The main theorem}
\label{sec:main}

\begin{theorem}
  \label{thm:main}
  \[
    v_{4}(9) = \frac{5}{6}\sin\frac{2\pi}{5}
      = \frac{5}{24}\sqrt{10 + 2\sqrt{5}} = 0.792547\ldots,
  \]
  and every $P \in \mathcal{P}_{4}(9)$ with $\vol_{4}(P) = v_{4}(9)$ is, up to
  an orthogonal transformation, the convex hull of a square and a regular
  pentagon inscribed in $\mathbb{S}^{3}$ in orthogonal $2$-planes through the
  origin.
\end{theorem}

\begin{proof}
  Let $V^{*} = \frac{1}{24}\cdot 4\sin\frac{2\pi}{4}\cdot 5\sin\frac{2\pi}{5} =
  \frac{5}{6}\sin\frac{2\pi}{5}$; since $\sin\frac{2\pi}{5} =
  \frac14\sqrt{10 + 2\sqrt5}$, this is the stated value. By
  Lemma~\ref{lem:freesum} and Corollary~\ref{cor:pair} with $(m, l) = (4, 5)$,
  the square--pentagon pair lies in $\mathcal{P}_{4}(9)$ and has volume
  $V^{*}$, so $v_{4}(9) \geq V^{*}$.

  Let $P \in \mathcal{P}_{4}(9)$ with $\vol_{4}(P) = v_{4}(9)$, which exists by
  Remark~\ref{rem:attained}. By Remark~\ref{rem:propertyz} and
  Lemma~\ref{lem:simplicial}, $P$ is simplicial. Since $39/50 < V^{*}$,
  Proposition~\ref{prop:rivals} shows that $P$ has the combinatorial type of
  \texttt{firsching-9-10}, which by Lemma~\ref{lem:data} is that of the
  polytope of Lemma~\ref{lem:freesum} with $(m, l) = (4, 5)$. Corollary~\ref{cor:pair} gives
  $v_{4}(9) = \vol_{4}(P) \leq V^{*}$, so $v_{4}(9) = V^{*}$, and its equality
  case identifies $P$.
\end{proof}
